\documentclass[10pt,a4paper]{amsart}
\usepackage[margin=19mm]{geometry}
\usepackage{amsmath,amssymb,mathtools}
\usepackage[scr=rsfs,cal=euler]{mathalfa}
\usepackage{xfrac}
\usepackage[unicode,hidelinks,bookmarks=false]{hyperref}
\newcommand{\ie}{i.e.\ }
\newcommand{\cf}{cf.\ }
\newcommand{\resp}{respectively\ }
\newcommand{\Subsection}{\subsection}
\providecommand{\nobreakdash}{\nobreak\hbox{-}}
\setlength{\emergencystretch}{2em}
\multlinegap0pt
\usepackage{amssymb}
\usepackage{xcolor}
\usepackage{algorithm}\usepackage{algpseudocode}
\usepackage{enumerate}
\usepackage{mathtools}
\newtheorem{rem}{Remark}
\newtheorem{theorem}{Theorem}
\newtheorem{lemma}{Lemma}
\newtheorem{proposition}{Proposition}
\DeclareMathOperator{\Per}{\textup{Per}}
\DeclareMathOperator{\diam}{\textup{diam}}
\DeclareMathOperator{\difm}{\textup{d$\mathcal{H}$}}
\DeclareMathOperator{\dist}{\textup{dist}}
\title{Existence of Anisotropic Minkowski Content}
\author{Filip Fry\v{s}}
\date{Source: arXiv:2508.08156v5 (CC BY 4.0)}
\begin{document}
\maketitle

\noindent\emph{Source and licence.} Existence of Anisotropic Minkowski Content. Version arXiv:2508.08156v5 (CC BY 4.0), distributed by arXiv under the Creative Commons Attribution 4.0 International licence, http://creativecommons.org/licenses/by/4.0/, as recorded in the arXiv licence metadata for this exact version and shown on the abstract page https://arxiv.org/abs/2508.08156v5. Full licence notice: \texttt{licenses/arxiv-2508.08156-CC-BY-4.0.txt}.\par\smallskip

\noindent\textit{Editorial scope.} Counted unit: the article's theorem theoremExistence (source0.tex lines 868--878). The article's complete original proof of it is delivered with it (source0.tex lines 879--983, one \texttt{proof} environment of 105 lines). No mathematics is written, repaired or re-derived: the statement and the proof are the article's own lines, in the article's own notation, and no retained line is substituted or re-flowed.  Retained uncounted as the source-local context the proof needs: lines 234--237; lines 333--341; lines 373--375; lines 480--514; lines 525--531; lines 533--545.  Nothing was omitted from any retained span, so the package carries no omission record.  The retained lines cite 3 works, whose entries are transcribed verbatim from the article's own bibliography source in the e-print and delivered with short editorial labels recorded in the item's reference mapping.  No retained line contains a figure environment, an image-inclusion command or drawing code, so no image is delivered and the package declares no asset dependency.  Editorial formatting only, in addition: an AMS \texttt{amsart} preamble that restates the article's own declarations and macros used by the retained lines, namely the environments \texttt{rem}, \texttt{theorem}, \texttt{lemma}, \texttt{proposition} on separate counters, together with the standard packages those lines need.  The retained environments print in this document as Remark 1, Theorem 1, Lemma 1, Proposition 1 in order of appearance, whereas the article prints the same items with the numbering of its own full document.  The complete native source of the article as distributed in the arXiv e-print for this exact version is bundled unchanged as \texttt{sources/183075/source0.tex}.
\begin{rem}\label{R1}
Let $S \subseteq \mathbb{R}^n$ be a countably $\mathcal{H}^{n-1}$-rectifiable set with $\mathcal{H}^{n-1}(S) < \infty$.  
Then there exists a countable family of pairwise disjoint compact subsets of $S$ that covers $S$ up to an $\mathcal{H}^{n-1}$-negligible set, each of which is contained in an $(n-1)$-dimensional Lipschitz graph and has finite $\mathcal{H}^{n-1}$-measure (see \cite[Lemma 11.1 and Remark 11.7]{simon}). 
\end{rem}

\begin{theorem}\label{bes}
Let $A \subseteq \mathbb{R}^n$ and $\rho \colon A \to (0, \infty)$ be a bounded function.  
Then there exist $N \in \mathbb{N}$, depending only on $n$, and an at most countable set $S \subseteq A$ such that
\begin{align*}
A &\subseteq \bigcup_{x \in S} B(x, \rho(x)), 
&
\sum_{x \in S} \chi_{B(x, \rho(x))} &\le N.
\end{align*}
\end{theorem}

    \begin{equation}\label{hurich}
        a\,h_C(\nu) \leq |\nu| \leq b\,h_C(\nu).
    \end{equation}

\begin{rem}\label{REMAk}
Suppose $E$ is a set of finite perimeter in an open set $\Omega\subseteq\mathbb{R}^n$, $C, C'\in\mathcal{C}^n_0$, and $S\subseteq\mathbb{R}^n$ is measurable. Then:
\begin{itemize}
    \item The $C$-anisotropic outer Minkowski content of $E$ in $\Omega$ is sometimes called one-sided. Intuitively, we enlarge $E$ only \emph{outwards}, whereas $\partial E$ is enlarged in both directions, \emph{inwards} and \emph{outwards}.
    \item The classical perimeter satisfies $\Per(E;\Omega) = \textup{Per}_{h_{B(0,1)}}(E;\Omega)$.
    \item If $a C\subseteq C'\subseteq b C$ for some positive constants $a$ and $b$, then
    \[
    a\,\textup{Per}_{h_C}(E;\Omega) \le \textup{Per}_{h_{C'}}(E;\Omega) \le b\,\textup{Per}_{h_C}(E;\Omega) \quad \text{(see (\ref{hurich}))},
    \]
    \[
    a\,\mathcal{M}_{a\varepsilon,C}(S;\Omega) \le \mathcal{M}_{\varepsilon,C'}(S;\Omega) \le b\,\mathcal{M}_{b\varepsilon,C}(S;\Omega), \quad \varepsilon>0,
    \]
    \[
    a\,\mathcal{SM}_{a\varepsilon,C}(E;\Omega) \le \mathcal{SM}_{\varepsilon,C'}(E;\Omega) \le b\,\mathcal{SM}_{b\varepsilon,C}(E;\Omega), \quad \varepsilon>0.
    \]
    Hence, these three functionals behave similarly. In particular, $\mathcal{SM}_C(E;\Omega)=0$ implies $\mathcal{SM}_{C'}(E;\Omega)=0$, and similarly $\mathcal{M}_C(S;\Omega)=0$ implies $\mathcal{M}_{C'}(S;\Omega)=0$.
    \item For $\varepsilon>0$ (see the discussion on page 5 in \cite{chambolle}):
    \begin{equation}\label{min4}
        \min_{\{F\subseteq \Omega : \lambda^n(E\Delta F)=0\}} \mathcal{SM}_{\varepsilon,C}(F;\Omega) = \mathcal{SM}_{\varepsilon,C}(E^1;\Omega) = \mathcal{SM}_{\varepsilon,C}(\Omega\setminus E^0;\Omega).
    \end{equation}
    \item For $\varepsilon>0$, the functionals $2\mathcal{M}_{\varepsilon,C}(\cdot;\Omega)$ and $\mathcal{SM}_{\varepsilon,C}(\cdot;\Omega)$ coincide on $\lambda^n$-negligible sets.
    \item For all $\varepsilon>0$, $\mathcal{M}_{\varepsilon,C}(S;\Omega) \ge \lambda^n(S\cap\Omega)/(2\varepsilon)$, and therefore if $\mathcal{M}_C(S;\Omega)^*<\infty$, then $\lambda^n(S\cap\Omega)=0$.
        \item For all $\varepsilon>0$,
    \[
    \mathcal{SM}_{\varepsilon, C}(E; \Omega)
    =\frac{1}{\varepsilon}\lambda^n\Big(\big((E\oplus \varepsilon C)\cap\Omega\big)\setminus E\Big)
    \geq
    \frac{1}{\varepsilon}\lambda^n\Big(\big(\overline{E}\cap\Omega\big)\setminus E\Big),
    \]
    because $\overline{E}\subseteq E\oplus\varepsilon C$. Consequently, if $\mathcal{SM}_{C}(E; \Omega)^*<\infty$, then
    \[
    \lambda^n\big((\overline{E}\cap\Omega)\setminus E\big)=0.
    \]
\end{itemize}
\end{rem}

\begin{lemma}\label{crucial}
    If $S$ is a compact set contained in a Lipschitz $(n-1)$-graph, then
    \[
        \mathcal{M}_{C}(S; \Omega)=\frac{1}{2}\int\limits_{S} \big(h_C(\nu_S)+h_C(-\nu_S)\big)\, \difm^{n-1}
    \]
    for any open set $\Omega$ with $S\subseteq\Omega\subseteq\mathbb{R}^n$, where $\nu_S\colon S\to\mathbb{S}^{n-1}$ is a unit normal to $S$ (see \cite[p. 3]{Villa2}). 
\end{lemma}

\begin{proposition}[Lower bounds]\label{eq:4}
Let $E$ be a set of finite perimeter in an open set $\Omega\subseteq\mathbb{R}^n$, and let $C\in\mathcal{C}^n_0$. Then
\begin{equation}
    \mathcal{SM}_{C}(E; \Omega)_*=\liminf_{\varepsilon\to 0_+} \mathcal{SM}_{\varepsilon,C}(E; \Omega) \geq \textup{Per}_{h_C}(E; \Omega),
\end{equation}
and
\begin{equation}
\begin{split}
    \mathcal{M}_{C}(\partial E; \Omega)_* &= \liminf_{\varepsilon\to 0_+} \mathcal{M}_{\varepsilon,C}(\partial E; \Omega) \geq \mathcal{M}_{C}(\partial^* E; \Omega)_* = \liminf_{\varepsilon\to 0_+} \mathcal{M}_{\varepsilon,C}(\partial^* E; \Omega)\\
    &\geq \frac{1}{2}\big(\textup{Per}_{h_C}(E; \Omega)+\textup{Per}_{h_C}(\Omega\setminus E; \Omega)\big).
\end{split}
\end{equation} 
\end{proposition}

\begin{theorem}
\label{theoremExistence}
Let $E$ be a set of finite perimeter in an open set $\Omega\subseteq\mathbb{R}^n$, and let $C, C^{\prime}\in\mathcal{C}^n_0$. Then
\[
\mathcal{M}_C(\partial E; \Omega)=\frac{1}{2}\big(\textup{Per}_{h_{C}}(E; \Omega)+\textup{Per}_{h_{C}}(\Omega\setminus E; \Omega)\big)
\]
if and only if
\[
\mathcal{M}_{C^{\prime}}(\partial E; \Omega)=\frac{1}{2}\big(\textup{Per}_{h_{C^{\prime}}}(E; \Omega)+\textup{Per}_{h_{C^{\prime}}}(\Omega\setminus E; \Omega)\big). 
\]
\end{theorem}

\begin{proof}
Since the roles of $C$ and $C^{\prime}$ are symmetric, it suffices to prove one implication.

Suppose that
\[
\mathcal{M}_C(\partial E; \Omega)=\frac{1}{2}\big(\textup{Per}_{h_{C}}(E; \Omega)+\textup{Per}_{h_{C}}(\Omega\setminus E; \Omega)\big).
\]
If $\mathcal{H}^{n-1}(\Omega\cap\partial^*E)=0$, the claim follows from Remark \ref{REMAk}. Otherwise, assume $\mathcal{H}^{n-1}(\Omega\cap\partial^*E)>0$. It suffices to show
\begin{align}\label{nikl}
\mathcal{M}_{C^{\prime}}(\partial E; \Omega)^*\leq\frac{1}{2}\big(\textup{Per}_{h_{C^{\prime}}}(E; \Omega)+\textup{Per}_{h_{C^{\prime}}}(\Omega\setminus E; \Omega)\big)\leq \mathcal{M}_{C^{\prime}}(\partial E; \Omega)_*.
\end{align}
The second inequality in (\ref{nikl}) always holds (see Proposition \ref{eq:4}), so it suffices to prove the first inequality. We decompose $\Omega\cap\partial^*E$ into finitely many Lipschitz pieces and a small remainder. The Lipschitz pieces are handled by Lemma \ref{crucial}, while the remainder is controlled by the Besicovitch covering theorem.

Since $\Omega\cap\partial^* E$ is countably $\mathcal{H}^{n-1}$-rectifiable and has finite $\mathcal{H}^{n-1}$-measure, there exists a sequence $\{S_k\}_{k=1}^{\infty}$ of pairwise disjoint compact subsets of $\Omega\cap\partial^* E$ that covers $\Omega\cap\partial^* E$ up to an $\mathcal{H}^{n-1}$-negligible set, each contained in the graph of a real-valued Lipschitz function (see Remark \ref{R1}). Since $\mathcal{H}^{n-1}(\Omega\cap\partial^* E)>0$, at least one of the sets $S_k$ is nonempty. Fix $k_0\in\mathbb{N}$ such that $S_{k_0}\neq\emptyset$. Then $\bigcup_{k=1}^K S_k\neq\emptyset$ for every $K\geq k_0$, so the distance function
\[
x\mapsto \dist\Bigl(x,\bigcup_{k=1}^K S_k\Bigr)
\]
is well defined.

We now prove the first inequality in (\ref{nikl}). Observe
\begin{equation*}
\begin{split}
\mathcal{M}_C(\partial E; \Omega)&=\frac{1}{2}\big(\textup{Per}_{h_{C}}(E; \Omega)+\textup{Per}_{h_{C}}(\Omega\setminus E; \Omega)\big)\\
&=\frac{1}{2}\int_{\partial^* E\cap\Omega}\big(h_C(\nu_E)+h_C(-\nu_E)\big)\,d\mathcal{H}^{n-1}\\
&=\sum_{k\in\mathbb{N}}\frac{1}{2}\int_{S_k}\big(h_C(\nu_E)+h_C(-\nu_E)\big)\,d\mathcal{H}^{n-1}=\sum_{k\in\mathbb{N}}\mathcal{M}_C(S_k; \Omega),
\end{split}
\end{equation*}
where the last equality follows from Lemma \ref{crucial}. For a given $\varepsilon>0$, there exists $K\geq k_0$ such that
\[
\mathcal{M}_C(\partial E; \Omega)-\sum_{k=1}^K\mathcal{M}_C(S_k; \Omega)<\varepsilon^{n}.
\]

Since $\{S_k\}_{k=1}^K$ is a finite family of pairwise disjoint compact subsets of $\Omega$, their mutual distances are positive. Therefore, for all sufficiently small $r>0$, the sets $S_k\oplus rC$, $k=1,\dots,K$, are pairwise disjoint. Hence
\begin{equation*}
\sum_{k=1}^K\mathcal{M}_C(S_k; \Omega)=\lim_{r\to0_+}\frac{1}{2r}\lambda^n\bigg(\Big(\bigcup_{k=1}^K S_k\oplus rC\Big)\cap\Omega\bigg),
\end{equation*}
so that
\begin{equation}\label{rum}
\lim_{r\to0_+}\frac{1}{2r}\lambda^n\Bigg(\Omega\cap\bigg(\Big((\Omega\cap\partial E)\oplus rC\Big)\setminus\Big(\bigcup_{k=1}^K S_k\oplus rC\Big)\bigg)\Bigg)<\varepsilon^n.
\end{equation}

Let positive constants $a$, $b$ and $c$ satisfy $B(0,a)\subseteq C$ and $B(0, c)\subseteq C^{\prime}\subseteq B(0,b)$, and define
\[
E_{r, \varepsilon}\coloneq \Big\{x\in \Omega\cap\partial E  : \dist\Big(x, \bigcup_{k=1}^K S_k\Big)>2\diam(C)r\varepsilon \Big\}\cap\Omega_{r, \varepsilon},
\]
where $\Omega_{r, \varepsilon}\coloneq \{x\in\Omega : \dist(x, \partial\Omega)>ar\varepsilon\}$ if $\partial\Omega\neq\emptyset$, and $\Omega_{r, \varepsilon}\coloneq \Omega$ otherwise. %There exists $r_0>0$ such that $E_{r, \varepsilon}\neq\emptyset$ for all $r\in(0,r_0)$. 

By Theorem \ref{bes}, there exists an at most countable set $I_{r, \varepsilon}\subseteq E_{r, \varepsilon}$ and $N\in\mathbb{N}$ such that
\begin{align}
    E_{r, \varepsilon}\subseteq \bigcup_{x\in I_{r, \varepsilon}} B(x, ar\varepsilon)&&\text{and}&&\sum_{x\in I_{r, \varepsilon}}\chi_{B(x, ar\varepsilon)}\leq N.
\end{align}
We have \begin{equation*}
    \begin{split}
        \mathcal{H}^0(I_{r, \varepsilon})\omega_n (ar\varepsilon)^n&=\sum_{x\in I_{r, \varepsilon}}\omega_n (ar\varepsilon)^n\\&=\sum_{x\in I_{r, \varepsilon}}\lambda^n\big(B(x, ar\varepsilon)\big)\\&\leq N\lambda^n\Big(\bigcup_{x\in I_{r, \varepsilon}}B(x, ar\varepsilon)\Big).
    \end{split}
\end{equation*} 
Notice
\[
\bigcup_{x\in I_{r, \varepsilon}} B(x, ar\varepsilon)\subseteq E_{r, \varepsilon}\oplus B(0, ar\varepsilon) \subseteq \bigg(\big((\Omega\cap\partial E)\oplus r\varepsilon C \big)\setminus \Big(\bigcup_{k=1}^K S_k\oplus r\varepsilon C\Big)\bigg)\cap\Omega.
\]
Since \begin{equation*}\begin{split}
\varepsilon^n&>\lim_{r\to0_+}\frac{1}{2r}\lambda^n\Bigg(\Omega\cap\bigg(\Big((\Omega\cap\partial E)\oplus rC\Big)\setminus\Big(\bigcup_{k=1}^K S_k\oplus rC\Big)\bigg)\Bigg)\\&=\lim_{r\to0_+}\frac{1}{2r\varepsilon}\lambda^n\Bigg(\Omega\cap\bigg(\Big((\Omega\cap\partial E)\oplus r\varepsilon C\Big)\setminus\Big(\bigcup_{k=1}^K S_k\oplus r\varepsilon C\Big)\bigg)\Bigg),
\end{split}
\end{equation*} inequality \eqref{rum} implies
\begin{equation}\label{frf}\limsup_{r\to0_+}\mathcal{H}^0(I_{r, \varepsilon}) r^{n-1}\leq \frac{2N\varepsilon}{a^n\omega_n}.
\end{equation}

We have $$\Omega\cap\partial E\subseteq \big((\Omega\cap\partial E)\setminus E_{r, \varepsilon}\big)\cup E_{r, \varepsilon}, \quad r>0.$$

On the one hand,
\[
(\partial E\cap\Omega)\setminus E_{r, \varepsilon}\subseteq\Big\{x\in \Omega\cap\partial E : \dist\Big(x, \bigcup_{k=1} ^K S_k\Big)\leq 2\diam(C)\varepsilon r\Big\}\cup\Omega^{\prime}_{r, \varepsilon}, \quad r>0,
\] where $$\Omega^{\prime}_{r, \varepsilon}\coloneq\{x\in \Omega\cap\partial E:\dist(x, \partial\Omega)\leq ar\varepsilon\}, \quad r>0,$$ provided that $\partial\Omega\neq\emptyset;$ otherwise, for $r>0$, put $\Omega^{\prime}_{r, \varepsilon}\coloneq\emptyset$.

Since $\Omega^{\prime}_{r, \varepsilon}\subseteq (\partial\Omega)\oplus B(0, ar\varepsilon)$ and $\bigcup_{k=1}^K S_k\subseteq\Omega$, we see that for any small enough $r>0$  $$(\Omega^{\prime}_{r, \varepsilon}\oplus rC^{\prime})\cap \Big(\bigcup_{k=1}^K S_k\oplus ba^{-1}rC\Big)=\emptyset,$$ whence $$(\Omega^{\prime}_{r, \varepsilon}\oplus rC^{\prime})\subseteq \big((\Omega\cap\partial E)\oplus ba^{-1}rC\big)\setminus\Big(\bigcup_{k=1}^K S_k\oplus ba^{-1}rC\Big),$$ which coupled with \eqref{rum} ensures that \begin{equation}\label{rumík}
    \limsup_{r\to0_+}\frac{1}{2r}\lambda^n\big((\Omega^{\prime}_{r, \varepsilon}\oplus rC^{\prime})\cap\Omega\big)<\varepsilon^n b^{-1}a
\end{equation} Furthermore, \begin{equation}\label{rumcajs}
    \begin{split}
        \lambda^n\bigg(\Big(\bigcup_{k=1}^K S_k\oplus B\big(0, 3\diam(C)\varepsilon r\big)\Big) \oplus rC^{\prime}\bigg)\leq \lambda^n\Big(\bigcup_{k=1}^K S_k\oplus r\big(1+c^{-1}3\diam(C)\varepsilon\big)C^{\prime}\Big).
    \end{split}
\end{equation}


On the other hand, 
\[
E_{r, \varepsilon} \oplus rC^{\prime}\subseteq \bigcup_{x\in I_{r, \varepsilon}} B(x, rb),
\]
whence, using \eqref{frf}, \begin{equation}\label{konec}
    \limsup_{r\to0_+}\frac{1}{2r}\lambda^n(E_{r, \varepsilon} \oplus rC^{\prime})\leq\frac{\omega_n b^n}{2}\limsup_{r\to0_+}\mathcal{H}^0(I_{r, \varepsilon})r^{n-1}\leq \frac{N b^n\varepsilon}{a^n}.\end{equation}

Finally, taking into account \eqref{rumík}, \eqref{rumcajs} and \eqref{konec}, we obtain
\[
\begin{split}
\mathcal{M}_{C^{\prime}}(\partial E; \Omega)^* &\leq \limsup_{r\to0_+}\frac{1}{2r}\lambda^n\bigg(\Big(\big((\Omega\cap \partial E)\setminus E_{r, \varepsilon}\big)\oplus rC^{\prime}\Big)\cap\Omega\bigg)\\&\quad\quad+\limsup_{r\to0_+}\frac{1}{2r}\lambda^n\big((E_{r, \varepsilon}\oplus rC^{\prime})\cap\Omega\big)\\
&\leq\limsup_{r\to0_+}\frac{1}{2r}\lambda^n\big((\Omega^{\prime}_{r, \varepsilon}\oplus rC^{\prime})\cap\Omega\big)\\&\quad\quad+\limsup_{r\to0_+}\frac{1}{2r}\lambda^n\Big(\bigcup_{k=1}^K S_k\oplus r\big(1+c^{-1}3\diam(C)\varepsilon\big)C^{\prime}\Big)\\&\quad\quad+\limsup_{r\to0_+}\frac{1}{2r}\lambda^n\big((E_{r, \varepsilon}\oplus rC^{\prime})\cap\Omega\big)\\&\leq \varepsilon^n b^{-1}a+\frac{Nb^n\varepsilon }{a^n} + \big(1+c^{-1}3\diam(C)\varepsilon\big)\sum_{k=1}^K\mathcal{M}_{C^{\prime}}(S_k; \Omega)\\&\leq \varepsilon^n b^{-1}a+\frac{Nb^n\varepsilon }{a^n} + \big(1+c^{-1}3\diam(C)\varepsilon\big)\sum_{k=1}^{\infty}\mathcal{M}_{C^{\prime}}(S_k; \Omega)\\
&= \varepsilon^n b^{-1}a+\frac{Nb^n\varepsilon }{a^n}\\&\quad\quad+ \big(1+c^{-1}3\diam(C)\varepsilon\big)\frac{1}{2}\big(\textup{Per}_{h_{C^{\prime}}}(E; \Omega)+\textup{Per}_{h_{C^{\prime}}}(\Omega\setminus E; \Omega)\big).
\end{split}
\]

Since $\varepsilon>0$ was arbitrary, we conclude
\[
\mathcal{M}_{C^{\prime}}(\partial E; \Omega)^* \leq \frac{1}{2}\big(\textup{Per}_{h_{C^{\prime}}}(E; \Omega)+\textup{Per}_{h_{C^{\prime}}}(\Omega\setminus E; \Omega)\big),
\]
which completes the proof.
\end{proof}

\begin{thebibliography}{99}
\bibitem[Villa2]{Villa2} Luca Lussardi and Elena Villa. A general formula for the anisotropic outer {M}inkowski content of a set. \textit{Proc. Roy. Soc. Edinburgh Sect. A}, 146:393--413, 2016.
\bibitem[chambo]{chambolle} Antonin Chambolle, Stefano Lisini and Luca Lussardi. A remark on the anisotropic outer {M}inkowski content. \textit{Advances in Calculus of Variations}, 7:241--266, 2014.
\bibitem[simon]{simon} Leon Simon. Lectures on Geometric Measure Theory. \textit{Centre for Mathematical Analysis, Australian National University}, 1984.
\end{thebibliography}
\end{document}
